\diarytitle{0049}{ヘヴィサイド関数を外力とする初期値問題のラプラス変換による解法}

単位階段関数 $u(t-a)$（$t<a$ で $0$、$t \ge a$ で $1$）は $\mathcal{L}[u(t-a)](s) = \dfrac{e^{-as}}{s}$ を満たし、第二移動定理 $\mathcal{L}[u(t-a)\,f(t-a)](s) = e^{-as} F(s)$ を逆向きに用いると $e^{-as}F(s)$ の逆ラプラス変換が $u(t-a)f(t-a)$ として求まる。

両辺をラプラス変換すると
\[
s^{2}Y(s) + 4Y(s) = \frac{e^{-s}}{s}
\quad\Longrightarrow\quad
Y(s) = e^{-s}\cdot\frac{1}{s(s^{2}+4)}
\]
まず $\dfrac{1}{s(s^{2}+4)}$ を部分分数分解する。$\dfrac{1}{s(s^{2}+4)} = \dfrac{A}{s} + \dfrac{Bs+C}{s^{2}+4}$ とおくと
$1 = A(s^{2}+4) + (Bs+C)s$。$s=0$ より $A = \dfrac{1}{4}$。$s^{2}$ の係数比較で $0 = A+B$ より $B=-\dfrac{1}{4}$。$s^{1}$ の係数比較で $C=0$。よって
\[
\frac{1}{s(s^{2}+4)} = \frac{1}{4}\cdot\frac{1}{s} - \frac{1}{4}\cdot\frac{s}{s^{2}+4}
\]
逆ラプラス変換すると
\[
\mathcal{L}^{-1}\!\left[\frac{1}{s(s^{2}+4)}\right] = \frac{1}{4}\left(1 - \cos 2t\right)
\]
第二移動定理より
\[
\boxed{\; y(t) = \frac{1}{4}\bigl(1 - \cos 2(t-1)\bigr)\,u(t-1) \;}
\]
検算: $t<1$ では $y=0$ で $y''+4y=0$（右辺も $0$）。$t>1$ では $y = \dfrac{1}{4}(1-\cos 2(t-1))$、$y' = \dfrac{1}{2}\sin 2(t-1)$、$y'' = \cos 2(t-1)$ であり、
\[
y''+4y = \cos 2(t-1) + \bigl(1-\cos 2(t-1)\bigr) = 1
\]
となり右辺 $u(t-1)=1$ と一致する。また $t=1$ で $y(1^{-})=0=y(1^{+})$、$y'(1^{-})=0=y'(1^{+})$ となり連続的につながる。
